\section{Artifact, split, and intervention details}
\label{app:artifacts}

\subsection{Pinned organism}

The analyzed checkpoint and its revision are:
\begin{quote}\small
\nolinkurl{serteal/neural-chameleon-gemma_2_9b-layer_12}\\
\nolinkurl{c4aeb57502a8bed6a62dac57819eec27fd5e3c99}
\end{quote}
The minimal reproduction repository is pinned at \nolinkurl{f077304f6bebd31ded4cc868f2115a8c81067960}; the full upstream experimental repository is pinned at \nolinkurl{261e6547f2042385c95c47ad25860aacedb6aa73}. Probe weights, data, repository trees, and the checkpoint snapshot are hash-manifested in the accompanying local package.

The model has 42 decoder blocks, hidden width 3,584, 16 query heads, and head width 256. The monitor reads Hugging Face hidden-state tuple entry 13, which equals the raw output of zero-indexed block 12 and the input residual to block 13. Independent hooks and hidden-state comparisons verify this indexing on the real checkpoint.

\subsection{Split accounting}

\begin{table}[h]
\centering
\caption{Frozen concept split. Counts are shown as positive + negative examples.}
\label{tab:splits}
\begin{tabular}{lll}
\toprule
Evidence role & Concepts & Per-concept counts \\
\midrule
Discovery & HTML, biology-focused, comforting, German & $64+64$ \\
Held-out benign & all-caps, chemistry-based, confused, Finnish, & $64+64$ \\
& jokey, literature-focused, mathematical & \\
Locked safety & deception & $62+62$ \\
Locked safety & harmfulness & $100+100$ \\
\bottomrule
\end{tabular}
\end{table}

Audits find no cross-phase example-ID, exact content, normalized prompt-response, source path/index, or word-five-gram Jaccard match at or above 0.8. This does not rule out paraphrastic similarity or unknown pretraining exposure.

\subsection{Frozen component population}
\label{app:components}

The K16 prefix was frozen in this discovery order:

\begin{enumerate}
  \item layer 11 MLP;
  \item layer 11 head 08;
  \item layer 10 head 12;
  \item layer 11 head 09;
  \item layer 12 MLP;
  \item layer 10 MLP;
  \item layer 10 head 02;
  \item layer 09 head 11;
  \item layer 12 head 12;
  \item layer 09 MLP;
  \item layer 12 head 03;
  \item layer 09 head 04;
  \item layer 11 head 14;
  \item layer 12 head 02;
  \item layer 09 head 13;
  \item layer 11 head 15.
\end{enumerate}

The ordering was chosen from exact discovery-concept rescue after screening a frozen 68-candidate universe. No identity or order changed after validation or safety access. Later decomposition shows that the 12 heads account for nearly all group effect on the Day 14 subset; the four whole-MLP sites remain listed because K16 was the confirmatory intervention actually frozen and tested.

\subsection{Behavior and confound controls}

Primary causal measurements target the activation probe, so behavior is measured separately. Fixed-continuation diagnostics include response NLL, next-token KL, and activation RMS. Prompt controls include monitoring-only, concept-only, irrelevant-concept, nearby-concept, paraphrased, malformed, and relocated triggers. Negative-example effects detect off-target probe motion. General-degradation and generation checks are secondary and cannot turn probe recovery into evidence of safety.

For the original K16 safety intervention, deception continuation behavior is comparatively stable (KL approximately 0.010 in both directions), while harmfulness is materially perturbed (KL 0.066 rescue and 0.054 induction; NLL shifts $-0.114$ and $+0.093$ nats/token). The later final-map equivalence result is narrower and must not be substituted for this K16 behavior result.
